% Options for packages loaded elsewhere
\PassOptionsToPackage{unicode}{hyperref}
\PassOptionsToPackage{hyphens}{url}
%
\documentclass[
]{article}
\usepackage{amsmath,amssymb}
\usepackage{iftex}
\ifPDFTeX
  \usepackage[T1]{fontenc}
  \usepackage[utf8]{inputenc}
  \usepackage{textcomp} % provide euro and other symbols
\else % if luatex or xetex
  \usepackage{unicode-math} % this also loads fontspec
  \defaultfontfeatures{Scale=MatchLowercase}
  \defaultfontfeatures[\rmfamily]{Ligatures=TeX,Scale=1}
\fi
\usepackage{lmodern}
\ifPDFTeX\else
  % xetex/luatex font selection
\fi
% Use upquote if available, for straight quotes in verbatim environments
\IfFileExists{upquote.sty}{\usepackage{upquote}}{}
\IfFileExists{microtype.sty}{% use microtype if available
  \usepackage[]{microtype}
  \UseMicrotypeSet[protrusion]{basicmath} % disable protrusion for tt fonts
}{}
\makeatletter
\@ifundefined{KOMAClassName}{% if non-KOMA class
  \IfFileExists{parskip.sty}{%
    \usepackage{parskip}
  }{% else
    \setlength{\parindent}{0pt}
    \setlength{\parskip}{6pt plus 2pt minus 1pt}}
}{% if KOMA class
  \KOMAoptions{parskip=half}}
\makeatother
\usepackage{xcolor}
\usepackage[margin=2.5cm]{geometry}
\usepackage{longtable,booktabs,array}
\usepackage{calc} % for calculating minipage widths
% Correct order of tables after \paragraph or \subparagraph
\usepackage{etoolbox}
\makeatletter
\patchcmd\longtable{\par}{\if@noskipsec\mbox{}\fi\par}{}{}
\makeatother
% Allow footnotes in longtable head/foot
\IfFileExists{footnotehyper.sty}{\usepackage{footnotehyper}}{\usepackage{footnote}}
\makesavenoteenv{longtable}
\setlength{\emergencystretch}{3em} % prevent overfull lines
\providecommand{\tightlist}{%
  \setlength{\itemsep}{0pt}\setlength{\parskip}{0pt}}
\setcounter{secnumdepth}{-\maxdimen} % remove section numbering
\ifLuaTeX
  \usepackage{selnolig}  % disable illegal ligatures
\fi
\IfFileExists{bookmark.sty}{\usepackage{bookmark}}{\usepackage{hyperref}}
\IfFileExists{xurl.sty}{\usepackage{xurl}}{} % add URL line breaks if available
\urlstyle{same}
\hypersetup{
  hidelinks,
  pdfcreator={LaTeX via pandoc}}

\author{}
\date{}

\begin{document}

\hypertarget{transforming-the-canonical-lagrangian-of-a-relativistic-particle-for-a-smooth-massless-limit-the-einbein-formulation-in-2-1-dimensions}{%
\section{Transforming the Canonical Lagrangian of a Relativistic
Particle for a Smooth Massless Limit: The Einbein Formulation in 2 + 1
Dimensions}\label{transforming-the-canonical-lagrangian-of-a-relativistic-particle-for-a-smooth-massless-limit-the-einbein-formulation-in-2-1-dimensions}}

Version 1.1 --- 11 August 2026

Giuliano (Vantasner AG), extending the 1 + 1 tutorial of Tiziano Fulceri
(DOI: 10.5281/zenodo.21879560)

Technical Tutorial Notes for Mathematical Physics Students (General
Relativity, Lagrangian Mechanics, and the Massless Limit)

\hypertarget{abstract}{%
\subsection{Abstract}\label{abstract}}

This tutorial extends, step by step and without omitting intermediate
algebraic manipulations, the einbein reformulation of the relativistic
point-particle action from 1 + 1 to 2 + 1 spacetime dimensions. The
variational core of the construction is dimension-independent and is
re-derived in full: all Euler--Lagrange equations are computed
explicitly, the algebraic solution for the einbein is substituted back
into the action, and the recovery of the original square-root form is
verified by direct calculation. The genuinely new content of the 2 +
1-dimensional setting is then isolated and developed: the null
directions form a circle rather than a pair of rays, the mass-shell
constraint underdetermines the velocity, the little groups are \(SO(2)\)
(massive) and \(\mathbb{R}\) (massless) rather than the trivial group,
and massive particles in 2 + 1 dimensions may carry any real spin
(anyons) while massless particles --- in the conventional physical
sector, where the translation-like massless little group is taken to act
trivially --- carry no spin degree of freedom at all (the mathematically
allowed alternative is discussed in Section 10.2). The presentation is
self-contained for a reader already familiar with the principles of
general relativity, the Lagrangian formalism, and the Euler--Lagrange
equations; familiarity with the companion 1 + 1-dimensional notes is
helpful but not assumed. References to standard textbooks and original
papers are provided with author names, titles, publication years and,
where available, stable URLs.

\hypertarget{contents}{%
\subsection{Contents}\label{contents}}

\begin{enumerate}
\def\labelenumi{\arabic{enumi}.}
\tightlist
\item
  Introduction and Motivation
\item
  Notation and Geometric Setup in 2 + 1 Dimensions
\item
  The Canonical Square-Root Action
\item
  Difficulties with the Massless Limit
\item
  The Einbein Reformulation
\item
  Euler--Lagrange Equation for the Einbein
\item
  Substitution and Recovery of the Square-Root Action
\item
  Euler--Lagrange Equations for the Coordinates (Massive Case)
\item
  The Massless Case: Euler--Lagrange Equations and Null Geodesics
\item
  Where the Dimension Enters: Null Directions, Mass Shells, and Little
  Groups
\item
  Brief Hamiltonian Perspective
\item
  Resolution of the Equations of Motion and the Resulting Trajectories
\item
  Summary
\end{enumerate}

\hypertarget{introduction-and-motivation}{%
\subsection{1 Introduction and
Motivation}\label{introduction-and-motivation}}

The companion tutorial {[}0{]} carried out the einbein reformulation of
the relativistic point particle in 1 + 1 spacetime dimensions, deriving
every algebraic step explicitly. The present notes perform the same
construction in 2 + 1 dimensions. The motivation for doing so is not
merely to add an index. In 1 + 1 dimensions the construction works, but
it never has to answer the questions that make the massless limit
physically interesting, because in 1 + 1 dimensions

\begin{itemize}
\tightlist
\item
  the null cone consists of exactly two rays, so a massless particle is
  characterized by a single discrete bit (left- or right-moving);
\item
  the null condition fixes a massless velocity up to its scale and that
  discrete choice, and even the massive mass shell retains only
  \textbf{one} continuous parameter (a rapidity), so the constraint
  analysis is almost trivial;
\item
  there is no transverse space, so no rotation group can act on the
  internal state of a particle, and spin in the ordinary sense does not
  exist.
\end{itemize}

In 2 + 1 dimensions each of these simplifications fails, in a controlled
and instructive way:

\begin{itemize}
\tightlist
\item
  the null directions form a \textbf{circle} \(S^{1}\) --- a continuous
  family --- rather than two points;
\item
  the mass-shell condition is one equation for two independent velocity
  ratios, so it \textbf{underdetermines} the velocity: the direction of
  motion is genuine initial data;
\item
  the massive little group is the rotation group \(SO(2)\), whose
  universal cover has irreducible unitary representations labeled by any
  real number, so massive particles in 2 + 1 dimensions can carry
  \textbf{anyonic spin}; the massless little group collapses to the
  translation group \(\mathbb{R}\), and in the conventional physical
  sector --- where this translation-like action is taken to be trivial
  --- massless states carry \textbf{no spin degree of freedom} at all
  (the nontrivial representations are mathematically allowed; see
  Section 10.2).
\end{itemize}

The einbein device itself, however, is indifferent to the dimension: the
same auxiliary world-line field \(e(\tau) > 0\) converts the square-root
action into a quadratic one, the same algebraic constraint appears, and
the same reparametrization gauge \(e = \mathrm{const}\) restores the
affine geodesic equation. The first half of these notes (Sections 2--9)
therefore re-derives the construction with all steps displayed, so that
the present document remains self-contained in the style of {[}0{]}. The
second half (Sections 10--12) develops the dimensional novelties listed
above, with references to the original literature.

A roadmap for the reader who has already studied {[}0{]}: Sections 2--9
are the 1 + 1 derivation with three indices instead of two, and can be
skimmed except for Section 8.4, where the eighteen independent
Christoffel components of a general 2 + 1 metric are written out.
Sections 10 and 12 contain the new material.

\hypertarget{notation-and-geometric-setup-in-2-1-dimensions}{%
\subsection{2 Notation and Geometric Setup in 2 + 1
Dimensions}\label{notation-and-geometric-setup-in-2-1-dimensions}}

We work on a three-dimensional Lorentzian manifold \(M\) equipped with a
metric tensor \(g_{\mu\nu}\) of signature \((-, +, +)\). Local
coordinates are denoted \(x^{\mu} = (x^{0}, x^{1}, x^{2})\). A
world-line is a smooth map

\[x^{\mu} : I \subset \mathbb{R} \longrightarrow M, \qquad \tau \longmapsto x^{\mu}(\tau), \tag{1}\]

where \(\tau\) is an arbitrary parameter (not necessarily proper time).
Derivatives with respect to \(\tau\) are written with a dot:
\(\dot{x}^{\mu}(\tau) := \dfrac{dx^{\mu}}{d\tau}\).

The induced interval on the world-line is

\[ds^{2} = g_{\mu\nu}(x(\tau))\, dx^{\mu} dx^{\nu} = g_{\mu\nu}(x)\, \dot{x}^{\mu} \dot{x}^{\nu}\, d\tau^{2}. \tag{2}\]

We adopt the convention that a future-directed timelike vector
\(v^{\mu}\) satisfies

\[g_{\mu\nu} v^{\mu} v^{\nu} < 0. \tag{3}\]

Consequently the quantity

\[\sigma := g_{\mu\nu}(x)\, \dot{x}^{\mu} \dot{x}^{\nu} \tag{4}\]

is negative for a timelike world-line. Because the signature is
\((-,+,+)\), the metric has one negative and two positive eigenvalues;
the determinant, being the product of the eigenvalues, is negative:

\[\det(g_{\mu\nu}) = g_{00}g_{11}g_{22} + 2\, g_{01}g_{12}g_{02} - g_{00}g_{12}^{2} - g_{11}g_{02}^{2} - g_{22}g_{01}^{2} < 0, \tag{5}\]

where the explicit expression is the Leibniz (cofactor) expansion of the
\(3 \times 3\) determinant: the three even permutations of the columns
contribute with a plus sign and the three odd permutations with a minus
sign, with the symmetry \(g_{\mu\nu} = g_{\nu\mu}\) used to collect the
two mixed products \(g_{01}g_{12}g_{20}\) and \(g_{02}g_{21}g_{10}\)
into the single term \(2\,g_{01}g_{12}g_{02}\).

Throughout we keep the speed of light \(c_{0}\) explicit. Rest mass is
denoted \(m \geq 0\).

\hypertarget{the-canonical-square-root-action}{%
\subsection{3 The Canonical Square-Root
Action}\label{the-canonical-square-root-action}}

The action that extremizes proper time is

\[S_{\mathrm{sqrt}}[x] = -mc_{0} \int_{\tau_{i}}^{\tau_{f}} \sqrt{-g_{\mu\nu}(x)\, \dot{x}^{\mu} \dot{x}^{\nu}}\; d\tau = -mc_{0} \int \sqrt{-\sigma}\; d\tau. \tag{6}\]

The associated Lagrangian density (with respect to the parameter
\(\tau\)) is therefore

\[L_{\mathrm{sqrt}} = -mc_{0}\sqrt{-g_{\mu\nu}\, \dot{x}^{\mu} \dot{x}^{\nu}} = -mc_{0}\sqrt{-\sigma}. \tag{7}\]

Because \(L_{\mathrm{sqrt}}\) is homogeneous of degree one in the
velocities, the action is invariant under arbitrary reparametrizations
\(\tau \to f(\tau)\) with \(\dot{f} > 0\).\footnote{``Reparametrization
  invariance'' means that the value of the action does not change if the
  world-line is relabeled by any strictly increasing function \(f\): the
  curve as a geometric object is unchanged, only its parametrization.
  Homogeneity of degree one,
  \(L(x, \lambda \dot{x}) = \lambda L(x, \dot{x})\) for \(\lambda > 0\),
  is precisely the property that makes the integral
  parameter-independent by the change-of-variables formula.} The
Euler--Lagrange equations obtained from \(L_{\mathrm{sqrt}}\) are
equivalent to the geodesic equation

\[\frac{D \dot{x}^{\mu}}{d\tau} = 0 \tag{8}\]

when the parameter is chosen to be an affine parameter\footnote{A
  parameter \(\tau\) along a geodesic is called affine if the tangent
  vector is parallel-transported without rescaling:
  \(D\dot{x}^{\mu}/d\tau = 0\). Any other parameter produces an extra
  term proportional to \(\dot{x}^{\mu}\); such a parameter is called
  non-affine. Affine parameters along a given geodesic are unique up to
  affine changes \(\tau \to a\tau + b\); see S. M. Carroll,
  \emph{Spacetime and Geometry} (Cambridge University Press, 2019),
  §3.3.} (or, more generally, proportional to proper time). None of
these statements involves the dimension of \(M\); the proofs given in
{[}0{]}, Section 3, apply verbatim.

\hypertarget{difficulties-with-the-massless-limit}{%
\subsection{4 Difficulties with the Massless
Limit}\label{difficulties-with-the-massless-limit}}

If one attempts to set \(m = 0\) directly in (7), the action vanishes
for every trajectory that satisfies the null condition \(\sigma = 0\).
The variational principle then ceases to select a preferred class of
curves: every null curve yields the same (zero) value of the action.
Moreover, the square-root Lagrangian becomes non-differentiable on the
light cone. Consequently a different formulation is required if one
wishes a single action principle that covers both the massive and the
massless cases continuously.

In 2 + 1 dimensions there is an additional, kinematic hint that the
massless theory must be approached with care. For a massive particle the
tangent vector is confined to the interior of the future light cone, a
connected open set in each tangent space. As \(m \to 0\) the tangent
vector is pushed onto the boundary of that cone. In 1 + 1 dimensions the
boundary of the future cone is a discrete set of two rays, and the
massive theory prepares the discrete label continuously: a timelike
velocity with \(\dot{x}^{1}/\dot{x}^{0} > 0\) tends to the right-moving
ray. In 2 + 1 dimensions the boundary of the future cone is a
two-dimensional cone, whose space of directions is a circle; the
continuous angular label survives the limit. Section 10 makes this
precise.

\hypertarget{the-einbein-reformulation}{%
\subsection{5 The Einbein
Reformulation}\label{the-einbein-reformulation}}

Introduce a real auxiliary field \(e(\tau) > 0\), called the einbein.
Geometrically, \(e\) may be viewed as defining an intrinsic metric on
the one-dimensional world-line via

\[ds^{2}_{\mathrm{worldline}} = e(\tau)^{2}\, d\tau^{2}. \tag{9}\]

Consider the new action functional that depends on both the embedding
\(x^{\mu}(\tau)\) and the einbein \(e(\tau)\):

\[S[x, e] = \int_{\tau_{i}}^{\tau_{f}} \left[ \frac{1}{2e}\, g_{\mu\nu}(x)\, \dot{x}^{\mu} \dot{x}^{\nu} - \frac{e}{2}(mc_{0})^{2} \right] d\tau. \tag{10}\]

The corresponding Lagrangian is

\[L(x, \dot{x}, e) = \frac{1}{2e}\, \sigma - \frac{e}{2}(mc_{0})^{2}, \tag{11}\]

where \(\sigma = g_{\mu\nu} \dot{x}^{\mu} \dot{x}^{\nu}\) as before.
Note that \(L\) is quadratic in the velocities and contains no
derivative of \(e\); the einbein is therefore an auxiliary
(non-dynamical) field.\footnote{An auxiliary (or non-dynamical) field is
  one whose equation of motion contains no derivatives of the field
  itself and can therefore be solved algebraically in terms of the other
  fields. Substituting the algebraic solution back into the action
  returns an equivalent description of the same classical theory.} The
action (10) is invariant under the reparametrizations

\[\tau \longmapsto f(\tau), \qquad e(\tau) \longmapsto \frac{e(\tau)}{\dot{f}(\tau)} \qquad (\dot{f} > 0), \tag{12}\]

because under the combined transformation the first term of the
integrand picks up \(1/\dot{f}^{2}\) from the two velocities and
\(\dot{f}\) from the measure, for a net \(1/\dot{f}\) --- which is
exactly cancelled by the einbein's law \(1/e \to \dot{f}/e\); the second
term picks up \(1/\dot{f}\) from \(e \to e/\dot{f}\) and \(\dot{f}\)
from the measure, and is invariant as it stands. The invariant
combination is \(e(\tau)\, d\tau\), the world-line volume element
\(\sqrt{g_{\mathrm{worldline}}}\, d\tau\) of (9). The einbein is thus
not an optional trick but the required compensating field: it is the
degree-one object on the world-line whose transformation law absorbs the
Jacobian \(\dot{f}\) so that a quadratic action can remain
reparametrization-invariant.

\hypertarget{eulerlagrange-equation-for-the-einbein}{%
\subsection{6 Euler--Lagrange Equation for the
Einbein}\label{eulerlagrange-equation-for-the-einbein}}

The general Euler--Lagrange equation associated with any variable
\(\varphi(\tau)\) is

\[\frac{d}{d\tau} \frac{\partial L}{\partial \dot{\varphi}} - \frac{\partial L}{\partial \varphi} = 0. \tag{13}\]

We now specialise this equation to the einbein by a sequence of purely
logical steps, exactly as in {[}0{]}, Section 6.

\textbf{Step 1.} Start from the general Euler--Lagrange equation (13).

\textbf{Step 2.} Choose the variable \(\varphi = e(\tau)\). The equation
becomes

\[\frac{d}{d\tau} \frac{\partial L}{\partial \dot{e}} - \frac{\partial L}{\partial e} = 0. \tag{14}\]

\textbf{Step 3.} Inspection of the explicit Lagrangian

\[L = \frac{1}{2e}\, g_{\mu\nu} \dot{x}^{\mu} \dot{x}^{\nu} - \frac{e}{2}(mc_{0})^{2} \tag{15}\]

shows that the derivative \(\dot{e}\) never appears. Consequently the
partial derivative with respect to that velocity vanishes identically:

\[\frac{\partial L}{\partial \dot{e}} \equiv 0. \tag{16}\]

\textbf{Step 4.} The ordinary derivative of the zero function is zero:

\[\frac{d}{d\tau} \frac{\partial L}{\partial \dot{e}} = \frac{d}{d\tau}(0) = 0. \tag{17}\]

\textbf{Step 5.} Substituting the result of Step 4 into the specialised
equation of Step 2 leaves only

\[- \frac{\partial L}{\partial e} = 0. \tag{18}\]

\textbf{Step 6.} Multiplying both sides by \(-1\) yields the equivalent
statement

\[\frac{\partial L}{\partial e} = 0. \tag{19}\]

Thus equation (19) is nothing but the general Euler--Lagrange equation
after the observation that \(\partial L / \partial \dot{e}\) is
identically zero has been used. No further dynamical assumption is
required; the reduction is purely formal.

We now compute the partial derivative \(\partial L / \partial e\) term
by term. The Lagrangian is the sum of two pieces:

\[L = \underbrace{\frac{1}{2e}\, \sigma}_{L_{1}} + \underbrace{\left( - \frac{e}{2}(mc_{0})^{2} \right)}_{L_{2}}.\]

Differentiating each piece separately with respect to \(e\) (treating
the velocity combination \(\sigma\) as independent of \(e\)) yields

\[\frac{\partial L_{1}}{\partial e} = \frac{\partial}{\partial e} \left( \frac{1}{2e} \sigma \right) = \sigma \cdot \frac{\partial}{\partial e}\left( \frac{1}{2e} \right) = \sigma \cdot \left( - \frac{1}{2e^{2}} \right) = - \frac{1}{2e^{2}}\, \sigma, \tag{20}\]

\[\frac{\partial L_{2}}{\partial e} = \frac{\partial}{\partial e} \left( - \frac{e}{2}(mc_{0})^{2} \right) = - \frac{1}{2}(mc_{0})^{2}. \tag{21}\]

Because differentiation is a linear operation,

\[\frac{\partial L}{\partial e} = \frac{\partial L_{1}}{\partial e} + \frac{\partial L_{2}}{\partial e} = - \frac{1}{2e^{2}}\, \sigma - \frac{1}{2}(mc_{0})^{2}. \tag{22}\]

The Euler--Lagrange condition \(\partial L / \partial e = 0\) (already
established in equation (19)) therefore requires

\[- \frac{1}{2e^{2}}\, \sigma - \frac{1}{2}(mc_{0})^{2} = 0. \tag{23}\]

To solve the algebraic equation we multiply both sides by the
nowhere-vanishing factor \(-2e^{2}\):

\[\left( -2e^{2} \right) \left( - \frac{1}{2e^{2}}\, \sigma \right) + \left( -2e^{2} \right) \left( - \frac{1}{2}(mc_{0})^{2} \right) = 0
\quad \Longrightarrow \quad \sigma + e^{2}(mc_{0})^{2} = 0. \tag{24}\]

Equivalently,

\[g_{\mu\nu}\, \dot{x}^{\mu} \dot{x}^{\nu} = - e^{2} (mc_{0})^{2}. \tag{25}\]

Solving for the positive root \(e > 0\) gives

\[e(\tau) = \frac{\sqrt{-\sigma}}{mc_{0}} = \frac{\sqrt{-g_{\mu\nu}(x)\, \dot{x}^{\mu} \dot{x}^{\nu}}}{mc_{0}}. \tag{26}\]

(The negative root would correspond to a time-reversed parametrization
and is discarded by the orientation convention \(e > 0\).)

Observe that the dimension of the spacetime has entered this section
only through the index range of \(\mu\) in \(\sigma\). Every step is
valid for any dimension, in particular for 2 + 1.

\hypertarget{substitution-and-recovery-of-the-square-root-action}{%
\subsection{7 Substitution and Recovery of the Square-Root
Action}\label{substitution-and-recovery-of-the-square-root-action}}

We now substitute the algebraic solution (26) back into the Lagrangian
(11) and verify that the original square-root Lagrangian is recovered
exactly.

From the constraint (25) we have

\[\sigma = - e^{2} (mc_{0})^{2}. \tag{27}\]

Insert this expression into the first term of \(L\):

\[\frac{1}{2e}\, \sigma = \frac{-e^{2}(mc_{0})^{2}}{2e} = - \frac{e}{2}(mc_{0})^{2}. \tag{28}\]

Therefore the whole Lagrangian becomes

\[L = - \frac{e}{2}(mc_{0})^{2} - \frac{e}{2}(mc_{0})^{2} = - e\, (mc_{0})^{2}. \tag{29}\]

Now replace \(e\) by its explicit solution (26):

\[L = - \left( \frac{\sqrt{-\sigma}}{mc_{0}} \right) (mc_{0})^{2} = - mc_{0} \sqrt{-\sigma} = - mc_{0} \sqrt{-g_{\mu\nu}\, \dot{x}^{\mu} \dot{x}^{\nu}}. \tag{30}\]

This is precisely \(L_{\mathrm{sqrt}}\) of equation (7). Consequently
the two action principles are classically equivalent for \(m > 0\):
every critical point of \(S[x, e]\) that satisfies the einbein equation
of motion projects to a critical point of \(S_{\mathrm{sqrt}}[x]\), and
vice versa. The equivalence holds in any spacetime dimension; the 2 +
1-dimensional case needed no modification.

\hypertarget{eulerlagrange-equations-for-the-coordinates-massive-case}{%
\subsection{8 Euler--Lagrange Equations for the Coordinates (Massive
Case)}\label{eulerlagrange-equations-for-the-coordinates-massive-case}}

We next derive the equations of motion for the embedding coordinates
\(x^{\lambda}(\tau)\) in the massive case \(m > 0\). The Lagrangian

\[L(x, \dot{x}, e) = \frac{1}{2e}\, g_{\mu\nu}(x)\, \dot{x}^{\mu} \dot{x}^{\nu} - \frac{e}{2}(mc_{0})^{2} \tag{31}\]

depends on the coordinates both through the metric coefficients
\(g_{\mu\nu}(x)\) and through the velocities \(\dot{x}^{\mu}\). The
Euler--Lagrange equation for each component \(x^{\lambda}\)
(\(\lambda = 0, 1, 2\)) is

\[\frac{d}{d\tau} \frac{\partial L}{\partial \dot{x}^{\lambda}} = \frac{\partial L}{\partial x^{\lambda}}. \tag{32}\]

Throughout this section the einbein \(e(\tau)\) is treated as an
independent field (its own equation of motion will be imposed
afterwards).

\hypertarget{partial-derivative-with-respect-to-the-velocity}{%
\subsubsection{8.1 Partial derivative with respect to the
velocity}\label{partial-derivative-with-respect-to-the-velocity}}

Differentiate \(L\) with respect to \(\dot{x}^{\lambda}\). Only the
first term contributes:

\[\frac{\partial L}{\partial \dot{x}^{\lambda}} = \frac{\partial}{\partial \dot{x}^{\lambda}} \left( \frac{1}{2e}\, g_{\mu\nu}(x)\, \dot{x}^{\mu} \dot{x}^{\nu} \right)
= \frac{1}{2e}\, g_{\mu\nu}(x) \left( \delta^{\mu}{}_{\lambda} \dot{x}^{\nu} + \dot{x}^{\mu} \delta^{\nu}{}_{\lambda} \right)
= \frac{1}{2e} \left( g_{\lambda\nu} \dot{x}^{\nu} + g_{\mu\lambda} \dot{x}^{\mu} \right)
= \frac{1}{e}\, g_{\lambda\nu}(x)\, \dot{x}^{\nu}, \tag{33}\]

where the last equality follows from the symmetry
\(g_{\mu\lambda} = g_{\lambda\mu}\). Consequently the left-hand side of
the Euler--Lagrange equation is the total \(\tau\)-derivative

\[\frac{d}{d\tau} \frac{\partial L}{\partial \dot{x}^{\lambda}} = \frac{d}{d\tau} \left( \frac{1}{e}\, g_{\lambda\nu}(x)\, \dot{x}^{\nu} \right). \tag{34}\]

\hypertarget{partial-derivative-with-respect-to-the-coordinate}{%
\subsubsection{8.2 Partial derivative with respect to the
coordinate}\label{partial-derivative-with-respect-to-the-coordinate}}

At fixed velocities the only dependence of \(L\) on the coordinate
\(x^{\lambda}\) comes from the metric tensor:

\[\frac{\partial L}{\partial x^{\lambda}} = \frac{1}{2e} \left( \partial_{\lambda} g_{\mu\nu}(x) \right) \dot{x}^{\mu} \dot{x}^{\nu}. \tag{35}\]

(The second term of \(L\) does not depend on \(x\).)

\hypertarget{the-eulerlagrange-equation-before-imposing-the-einbein-constraint}{%
\subsubsection{8.3 The Euler--Lagrange equation before imposing the
einbein
constraint}\label{the-eulerlagrange-equation-before-imposing-the-einbein-constraint}}

Equating (34) and (35) yields the explicit Euler--Lagrange equation

\[\frac{d}{d\tau} \left( \frac{1}{e}\, g_{\lambda\nu} \dot{x}^{\nu} \right) = \frac{1}{2e} \left( \partial_{\lambda} g_{\mu\nu} \right) \dot{x}^{\mu} \dot{x}^{\nu}. \tag{36}\]

We now convert (36) into an equivalent form by a sequence of elementary
calculus steps.

\textbf{Step 1.} The left-hand side of (36) is the ordinary derivative
of a product of two \(\tau\)-dependent factors, \(\dfrac{1}{e(\tau)}\)
and \(g_{\lambda\nu}(x(\tau))\, \dot{x}^{\nu}(\tau)\). Apply the
ordinary product rule:

\[\frac{d}{d\tau} \left( \frac{1}{e}\, g_{\lambda\nu} \dot{x}^{\nu} \right) = \left( \frac{d}{d\tau} \frac{1}{e} \right) \left( g_{\lambda\nu} \dot{x}^{\nu} \right) + \frac{1}{e}\, \frac{d}{d\tau} \left( g_{\lambda\nu} \dot{x}^{\nu} \right). \tag{37}\]

\textbf{Step 2.} Differentiate the reciprocal of the einbein by the
chain rule:

\[\frac{d}{d\tau} \frac{1}{e} = \frac{d}{d\tau} \left( e^{-1} \right) = - e^{-2}\, \dot{e} = - \frac{\dot{e}}{e^{2}}. \tag{38}\]

\textbf{Step 3.} Substitute the result of Step 2 back into the product
rule of Step 1:

\[\frac{d}{d\tau} \left( \frac{1}{e}\, g_{\lambda\nu} \dot{x}^{\nu} \right) = - \frac{\dot{e}}{e^{2}}\, \left( g_{\lambda\nu} \dot{x}^{\nu} \right) + \frac{1}{e}\, \frac{d}{d\tau} \left( g_{\lambda\nu} \dot{x}^{\nu} \right). \tag{39}\]

\textbf{Step 4.} Equation (36) therefore reads

\[- \frac{\dot{e}}{e^{2}}\, g_{\lambda\nu} \dot{x}^{\nu} + \frac{1}{e}\, \frac{d}{d\tau} \left( g_{\lambda\nu} \dot{x}^{\nu} \right) = \frac{1}{2e} \left( \partial_{\lambda} g_{\mu\nu} \right) \dot{x}^{\mu} \dot{x}^{\nu}. \tag{40}\]

\textbf{Step 5.} Multiply every term of the equation by the positive
factor \(e\) (which never vanishes by definition of the einbein). The
first term becomes
\(- \dfrac{\dot{e}}{e}\, g_{\lambda\nu} \dot{x}^{\nu}\), the second term
becomes
\(\dfrac{d}{d\tau} \left( g_{\lambda\nu} \dot{x}^{\nu} \right)\), and
the right-hand side becomes
\(\dfrac{1}{2} \left( \partial_{\lambda} g_{\mu\nu} \right) \dot{x}^{\mu} \dot{x}^{\nu}\).

\textbf{Step 6.} Rearrange the resulting terms to obtain

\[\frac{d}{d\tau} \left( g_{\lambda\nu} \dot{x}^{\nu} \right) - \frac{\dot{e}}{e}\, g_{\lambda\nu} \dot{x}^{\nu} = \frac{1}{2} \left( \partial_{\lambda} g_{\mu\nu} \right) \dot{x}^{\mu} \dot{x}^{\nu}. \tag{41}\]

All six steps are ordinary calculus (product rule and chain rule)
together with multiplication by a nowhere-vanishing positive function;
no geometric identity is required for this particular transition, and no
step refers to the spacetime dimension.

\hypertarget{expansion-of-the-total-derivative-and-introduction-of-the-christoffel-symbols}{%
\subsubsection{8.4 Expansion of the total derivative and introduction of
the Christoffel
symbols}\label{expansion-of-the-total-derivative-and-introduction-of-the-christoffel-symbols}}

Expand the total derivative that appears on the left-hand side of (41).
Write the product as \(g_{\lambda\nu}(x(\tau))\, \dot{x}^{\nu}(\tau)\)
and apply the ordinary product rule of one-variable calculus:

\[\frac{d}{d\tau} \left( g_{\lambda\nu} \dot{x}^{\nu} \right) = \left( \frac{d}{d\tau} g_{\lambda\nu} \right) \dot{x}^{\nu} + g_{\lambda\nu}\, \ddot{x}^{\nu}. \tag{42}\]

The second term is immediate. The first term requires care because the
metric components are functions of the spacetime coordinates, which
themselves depend on \(\tau\). In 2 + 1 dimensions this dependence may
be written explicitly:

\[g_{\lambda\nu} = g_{\lambda\nu}\left( x^{0}(\tau), x^{1}(\tau), x^{2}(\tau) \right).\]

The ordinary chain rule of multivariable calculus then states

\[\frac{d}{d\tau} g_{\lambda\nu} = \frac{\partial g_{\lambda\nu}}{\partial x^{0}}\, \dot{x}^{0} + \frac{\partial g_{\lambda\nu}}{\partial x^{1}}\, \dot{x}^{1} + \frac{\partial g_{\lambda\nu}}{\partial x^{2}}\, \dot{x}^{2}. \tag{43}\]

In index notation the sum over the three coordinates is abbreviated

\[\frac{d}{d\tau} g_{\lambda\nu} = \left( \partial_{\rho} g_{\lambda\nu} \right) \dot{x}^{\rho}, \tag{44}\]

where \(\partial_{\rho} = \partial / \partial x^{\rho}\) and the
repeated index \(\rho\) is summed over \(0, 1, 2\). Multiplying by the
remaining velocity \(\dot{x}^{\nu}\) produces the quadratic term

\[\left( \frac{d}{d\tau} g_{\lambda\nu} \right) \dot{x}^{\nu} = \left( \partial_{\rho} g_{\lambda\nu} \right) \dot{x}^{\rho} \dot{x}^{\nu}. \tag{45}\]

Collecting both contributions yields the expansion

\[\frac{d}{d\tau} \left( g_{\lambda\nu} \dot{x}^{\nu} \right) = \left( \partial_{\rho} g_{\lambda\nu} \right) \dot{x}^{\rho} \dot{x}^{\nu} + g_{\lambda\nu}\, \ddot{x}^{\nu}. \tag{46}\]

Substitution of this expansion into (41) produces

\[\left( \partial_{\rho} g_{\lambda\nu} \right) \dot{x}^{\rho} \dot{x}^{\nu} + g_{\lambda\nu}\, \ddot{x}^{\nu} - \frac{\dot{e}}{e}\, g_{\lambda\nu} \dot{x}^{\nu} = \frac{1}{2} \left( \partial_{\lambda} g_{\mu\nu} \right) \dot{x}^{\mu} \dot{x}^{\nu}. \tag{47}\]

We now raise the free (covariant) index \(\lambda\) by a sequence of
purely algebraic steps that use only the inverse metric.

\textbf{Step 1.} Move every term of (47) to the left-hand side:

\[g_{\lambda\nu}\, \ddot{x}^{\nu} + \left( \partial_{\rho} g_{\lambda\nu} \right) \dot{x}^{\rho} \dot{x}^{\nu} - \frac{1}{2} \left( \partial_{\lambda} g_{\mu\nu} \right) \dot{x}^{\mu} \dot{x}^{\nu} - \frac{\dot{e}}{e}\, g_{\lambda\nu} \dot{x}^{\nu} = 0. \tag{48}\]

\textbf{Step 2.} Multiply the entire equation by the inverse-metric
component \(g^{\sigma\lambda}\) and sum over \(\lambda\):

\[g^{\sigma\lambda} \left( g_{\lambda\nu}\, \ddot{x}^{\nu} + \left( \partial_{\rho} g_{\lambda\nu} \right) \dot{x}^{\rho} \dot{x}^{\nu} - \frac{1}{2} \left( \partial_{\lambda} g_{\mu\nu} \right) \dot{x}^{\mu} \dot{x}^{\nu} - \frac{\dot{e}}{e}\, g_{\lambda\nu} \dot{x}^{\nu} \right) = 0. \tag{49}\]

\textbf{Step 3.} Distribute the factor \(g^{\sigma\lambda}\) term by
term.

\begin{itemize}
\tightlist
\item
  First term (second derivatives):
  \(g^{\sigma\lambda} g_{\lambda\nu}\, \ddot{x}^{\nu} = \delta^{\sigma}{}_{\nu}\, \ddot{x}^{\nu} = \ddot{x}^{\sigma}\)
  (by the defining property of the inverse metric).
\item
  Last term (einbein derivative):
  \(g^{\sigma\lambda} \left( - \dfrac{\dot{e}}{e}\, g_{\lambda\nu} \dot{x}^{\nu} \right) = - \dfrac{\dot{e}}{e}\, \delta^{\sigma}{}_{\nu}\, \dot{x}^{\nu} = - \dfrac{\dot{e}}{e}\, \dot{x}^{\sigma}\).
\item
  The two middle terms that involve derivatives of the metric remain
  \(g^{\sigma\lambda} \left( \left( \partial_{\rho} g_{\lambda\nu} \right) \dot{x}^{\rho} \dot{x}^{\nu} - \frac{1}{2} \left( \partial_{\lambda} g_{\mu\nu} \right) \dot{x}^{\mu} \dot{x}^{\nu} \right)\).
\end{itemize}

\textbf{Step 4.} Because the velocity bilinears are symmetric, the
combination of metric derivatives that multiplies
\(\dot{x}^{\rho} \dot{x}^{\nu}\) may be rewritten

\[\frac{1}{2}\, g^{\sigma\lambda} \left( 2\, \partial_{\rho} g_{\lambda\nu} - \partial_{\lambda} g_{\rho\nu} \right) \dot{x}^{\rho} \dot{x}^{\nu}. \tag{50}\]

The two index placements are equal because renaming the summed indices
\(\rho \leftrightarrow \nu\) maps
\(\partial_{\nu} g_{\lambda\rho}\, \dot{x}^{\rho} \dot{x}^{\nu}\) onto
\(\partial_{\rho} g_{\lambda\nu}\, \dot{x}^{\rho} \dot{x}^{\nu}\); the
symmetrized combination is exactly what the Christoffel symbol requires.

\textbf{Step 5.} Collecting all pieces yields

\[\ddot{x}^{\sigma} + \frac{1}{2}\, g^{\sigma\lambda} \left( \partial_{\rho} g_{\lambda\nu} + \partial_{\nu} g_{\lambda\rho} - \partial_{\lambda} g_{\rho\nu} \right) \dot{x}^{\rho} \dot{x}^{\nu} - \frac{\dot{e}}{e}\, \dot{x}^{\sigma} = 0. \tag{51}\]

Every step above is ordinary linear algebra with the inverse metric
together with the symmetry of the velocity bilinears; no geometric
identity beyond the existence of the inverse metric is required.

The combination of metric derivatives in (51) is the Christoffel symbol
of the Levi-Civita connection,

\[\Gamma^{\sigma}{}_{\rho\nu} = \frac{1}{2}\, g^{\sigma\lambda} \left( \partial_{\rho} g_{\lambda\nu} + \partial_{\nu} g_{\lambda\rho} - \partial_{\lambda} g_{\rho\nu} \right). \tag{52}\]

Because the lower indices are symmetric,
\(\Gamma^{\sigma}{}_{\rho\nu} = \Gamma^{\sigma}{}_{\nu\rho}\), the
number of independent components in \(n\) spacetime dimensions is \(n\)
(choices of \(\sigma\)) times \(n(n+1)/2\) (symmetric pairs
\((\rho, \nu)\)), i.e.~\(n^{2}(n+1)/2\): six components in 1 + 1
dimensions, \textbf{eighteen in 2 + 1 dimensions}. For completeness, and
because the component bookkeeping is where higher-dimensional explicit
calculations most often go wrong, we list all eighteen, organized by the
upper index. For \(\sigma = 0\):

\[\begin{aligned}
\Gamma^{0}{}_{00} &= \tfrac{1}{2} g^{00}\, \partial_{0} g_{00} + \tfrac{1}{2} g^{01} \left( 2 \partial_{0} g_{10} - \partial_{1} g_{00} \right) + \tfrac{1}{2} g^{02} \left( 2 \partial_{0} g_{20} - \partial_{2} g_{00} \right), \\
\Gamma^{0}{}_{01} &= \tfrac{1}{2} g^{00}\, \partial_{1} g_{00} + \tfrac{1}{2} g^{01}\, \partial_{0} g_{11} + \tfrac{1}{2} g^{02} \left( \partial_{0} g_{12} + \partial_{1} g_{02} - \partial_{2} g_{01} \right), \\
\Gamma^{0}{}_{02} &= \tfrac{1}{2} g^{00}\, \partial_{2} g_{00} + \tfrac{1}{2} g^{01} \left( \partial_{0} g_{12} + \partial_{2} g_{01} - \partial_{1} g_{02} \right) + \tfrac{1}{2} g^{02}\, \partial_{0} g_{22}, \\
\Gamma^{0}{}_{11} &= \tfrac{1}{2} g^{00} \left( 2 \partial_{1} g_{01} - \partial_{0} g_{11} \right) + \tfrac{1}{2} g^{01}\, \partial_{1} g_{11} + \tfrac{1}{2} g^{02} \left( 2 \partial_{1} g_{21} - \partial_{2} g_{11} \right), \\
\Gamma^{0}{}_{12} &= \tfrac{1}{2} g^{00} \left( \partial_{1} g_{02} + \partial_{2} g_{01} - \partial_{0} g_{12} \right) + \tfrac{1}{2} g^{01}\, \partial_{2} g_{11} + \tfrac{1}{2} g^{02}\, \partial_{1} g_{22}, \\
\Gamma^{0}{}_{22} &= \tfrac{1}{2} g^{00} \left( 2 \partial_{2} g_{02} - \partial_{0} g_{22} \right) + \tfrac{1}{2} g^{01} \left( 2 \partial_{2} g_{12} - \partial_{1} g_{22} \right) + \tfrac{1}{2} g^{02}\, \partial_{2} g_{22}.
\end{aligned} \tag{53}\]

For \(\sigma = 1\):

\[\begin{aligned}
\Gamma^{1}{}_{00} &= \tfrac{1}{2} g^{10}\, \partial_{0} g_{00} + \tfrac{1}{2} g^{11} \left( 2 \partial_{0} g_{10} - \partial_{1} g_{00} \right) + \tfrac{1}{2} g^{12} \left( 2 \partial_{0} g_{20} - \partial_{2} g_{00} \right), \\
\Gamma^{1}{}_{01} &= \tfrac{1}{2} g^{10}\, \partial_{1} g_{00} + \tfrac{1}{2} g^{11}\, \partial_{0} g_{11} + \tfrac{1}{2} g^{12} \left( \partial_{0} g_{12} + \partial_{1} g_{02} - \partial_{2} g_{01} \right), \\
\Gamma^{1}{}_{02} &= \tfrac{1}{2} g^{10}\, \partial_{2} g_{00} + \tfrac{1}{2} g^{11} \left( \partial_{0} g_{12} + \partial_{2} g_{01} - \partial_{1} g_{02} \right) + \tfrac{1}{2} g^{12}\, \partial_{0} g_{22}, \\
\Gamma^{1}{}_{11} &= \tfrac{1}{2} g^{10} \left( 2 \partial_{1} g_{01} - \partial_{0} g_{11} \right) + \tfrac{1}{2} g^{11}\, \partial_{1} g_{11} + \tfrac{1}{2} g^{12} \left( 2 \partial_{1} g_{21} - \partial_{2} g_{11} \right), \\
\Gamma^{1}{}_{12} &= \tfrac{1}{2} g^{10} \left( \partial_{1} g_{02} + \partial_{2} g_{01} - \partial_{0} g_{12} \right) + \tfrac{1}{2} g^{11}\, \partial_{2} g_{11} + \tfrac{1}{2} g^{12}\, \partial_{1} g_{22}, \\
\Gamma^{1}{}_{22} &= \tfrac{1}{2} g^{10} \left( 2 \partial_{2} g_{02} - \partial_{0} g_{22} \right) + \tfrac{1}{2} g^{11} \left( 2 \partial_{2} g_{12} - \partial_{1} g_{22} \right) + \tfrac{1}{2} g^{12}\, \partial_{2} g_{22}.
\end{aligned} \tag{54}\]

For \(\sigma = 2\):

\[\begin{aligned}
\Gamma^{2}{}_{00} &= \tfrac{1}{2} g^{20}\, \partial_{0} g_{00} + \tfrac{1}{2} g^{21} \left( 2 \partial_{0} g_{10} - \partial_{1} g_{00} \right) + \tfrac{1}{2} g^{22} \left( 2 \partial_{0} g_{20} - \partial_{2} g_{00} \right), \\
\Gamma^{2}{}_{01} &= \tfrac{1}{2} g^{20}\, \partial_{1} g_{00} + \tfrac{1}{2} g^{21}\, \partial_{0} g_{11} + \tfrac{1}{2} g^{22} \left( \partial_{0} g_{12} + \partial_{1} g_{02} - \partial_{2} g_{01} \right), \\
\Gamma^{2}{}_{02} &= \tfrac{1}{2} g^{20}\, \partial_{2} g_{00} + \tfrac{1}{2} g^{21} \left( \partial_{0} g_{12} + \partial_{2} g_{01} - \partial_{1} g_{02} \right) + \tfrac{1}{2} g^{22}\, \partial_{0} g_{22}, \\
\Gamma^{2}{}_{11} &= \tfrac{1}{2} g^{20} \left( 2 \partial_{1} g_{01} - \partial_{0} g_{11} \right) + \tfrac{1}{2} g^{21}\, \partial_{1} g_{11} + \tfrac{1}{2} g^{22} \left( 2 \partial_{1} g_{21} - \partial_{2} g_{11} \right), \\
\Gamma^{2}{}_{12} &= \tfrac{1}{2} g^{20} \left( \partial_{1} g_{02} + \partial_{2} g_{01} - \partial_{0} g_{12} \right) + \tfrac{1}{2} g^{21}\, \partial_{2} g_{11} + \tfrac{1}{2} g^{22}\, \partial_{1} g_{22}, \\
\Gamma^{2}{}_{22} &= \tfrac{1}{2} g^{20} \left( 2 \partial_{2} g_{02} - \partial_{0} g_{22} \right) + \tfrac{1}{2} g^{21} \left( 2 \partial_{2} g_{12} - \partial_{1} g_{22} \right) + \tfrac{1}{2} g^{22}\, \partial_{2} g_{22}.
\end{aligned} \tag{55}\]

Each line is equation (52) with the sum over \(\lambda\) written out and
the trivial cancellations carried through (for instance, in
\(\Gamma^{\sigma}{}_{01}\) the \(\lambda = 0\) contribution
\(\partial_{0}g_{01} + \partial_{1}g_{00} - \partial_{0}g_{01}\)
collapses to \(\partial_{1}g_{00}\), and the \(\lambda = 1\)
contribution
\(\partial_{0}g_{11} + \partial_{1}g_{10} - \partial_{1}g_{01}\)
collapses to \(\partial_{0}g_{11}\)). The remaining nine components with
\(\rho > \nu\) (the pairs \((1,0)\), \((2,0)\), \((2,1)\) for each of
the three values of \(\sigma\)) are fixed by the symmetry
\(\Gamma^{\sigma}{}_{\rho\nu} = \Gamma^{\sigma}{}_{\nu\rho}\).

With these symbols the Euler--Lagrange equation takes the compact form

\[\ddot{x}^{\sigma} + \Gamma^{\sigma}{}_{\rho\nu}\, \dot{x}^{\rho} \dot{x}^{\nu} - \frac{\dot{e}}{e}\, \dot{x}^{\sigma} = 0, \tag{56}\]

or, equivalently,

\[\frac{D}{d\tau} \left( \frac{\dot{x}^{\sigma}}{e} \right) = 0. \tag{57}\]

Here the symbol \(\dfrac{D}{d\tau}\) denotes the covariant (or absolute)
derivative along the world-line. By definition,

\[\frac{D V^{\sigma}}{d\tau} := \frac{d V^{\sigma}}{d\tau} + \Gamma^{\sigma}{}_{\rho\nu}\, \dot{x}^{\rho} V^{\nu} \tag{58}\]

for any vector field \(V^{\sigma}\) defined along the curve.

Equation (56) is the geodesic equation written with respect to a
non-affine parameter. The term proportional to \(\dot{e}/e\) does not
disappear automatically; it is the precise measure of the failure of
\(\tau\) to be an affine parameter. The resolution is identical to the 1
+ 1-dimensional case, and we record it in the same five-step form:

\textbf{Step 1 --- Residual reparametrization freedom.} The pair
\((x(\tau), e(\tau))\) is defined only up to the transformations (12).
This residual freedom has not yet been fixed.

\textbf{Step 2 --- The algebraic (einbein) constraint.} Independently of
the choice of parameter, the einbein equation of motion is the purely
algebraic relation (25).

\textbf{Step 3 --- Imposing a convenient gauge condition.} Because of
the residual freedom one is allowed to impose one additional condition
on the variables. The most convenient choice is the gauge
condition\footnote{In the language of constrained systems a gauge
  condition is an extra equation that one is free to impose in order to
  remove the redundancy associated with a local symmetry. It is not an
  equation of motion; it is a choice of representative inside each
  equivalence class of reparametrizations.}

\[e(\tau) = \mathrm{const} \quad \Longrightarrow \quad \dot{e} = 0. \tag{59}\]

(This gauge is always attainable: given any positive function
\(e(\tau)\) one can solve the ordinary differential equation that
defines the new parameter \(f\) so that the transformed einbein is
constant.)

\textbf{Step 4 --- Disappearance of the non-affine term.} Once the gauge
condition of Step 3 is imposed, the coefficient of the unwanted term
vanishes identically, and equation (56) collapses to the canonical
affine geodesic equation

\[\frac{D \dot{x}^{\sigma}}{d\tau} = \ddot{x}^{\sigma} + \Gamma^{\sigma}{}_{\rho\nu}\, \dot{x}^{\rho} \dot{x}^{\nu} = 0. \tag{60}\]

\textbf{Step 5 --- Interpretation after the gauge choice.} With \(e\)
now a constant, the algebraic constraint (25) becomes a simple
normalisation of the tangent vector,

\[g_{\mu\nu}\, \dot{x}^{\mu} \dot{x}^{\nu} = - e^{2} (mc_{0})^{2} = \mathrm{const} < 0, \tag{61}\]

which identifies the parameter \(\tau\) with a multiple of the proper
time.

Thus the Euler--Lagrange equations of the einbein action, after the
einbein constraint is taken into account and a convenient gauge is
chosen, are exactly the geodesic equations of the background metric ---
in 2 + 1 dimensions exactly as in 1 + 1.

\hypertarget{the-massless-case-eulerlagrange-equations-and-null-geodesics}{%
\subsection{9 The Massless Case: Euler--Lagrange Equations and Null
Geodesics}\label{the-massless-case-eulerlagrange-equations-and-null-geodesics}}

We now set the rest mass to zero, \(m = 0\). The Lagrangian simplifies
to

\[L = \frac{1}{2e}\, g_{\mu\nu}(x)\, \dot{x}^{\mu} \dot{x}^{\nu}. \tag{62}\]

\hypertarget{eulerlagrange-equation-for-the-einbein-massless}{%
\subsubsection{9.1 Euler--Lagrange equation for the einbein
(massless)}\label{eulerlagrange-equation-for-the-einbein-massless}}

Because \(L\) still does not depend on \(\dot{e}\), the equation
\(\partial L / \partial e = 0\) reduces to

\[- \frac{1}{2e^{2}}\, g_{\mu\nu}\, \dot{x}^{\mu} \dot{x}^{\nu} = 0, \tag{63}\]

which immediately implies the null constraint

\[g_{\mu\nu}\, \dot{x}^{\mu} \dot{x}^{\nu} = 0. \tag{64}\]

\hypertarget{eulerlagrange-equations-for-the-coordinates-massless}{%
\subsubsection{9.2 Euler--Lagrange equations for the coordinates
(massless)}\label{eulerlagrange-equations-for-the-coordinates-massless}}

The partial derivatives with respect to velocity and coordinate are
identical in form to those of the massive case (the mass term is absent,
but it never contributed to these derivatives). Consequently one obtains
again (36), and exactly the same algebraic steps that led from (36) to
(56) now produce

\[\ddot{x}^{\sigma} + \Gamma^{\sigma}{}_{\rho\nu}\, \dot{x}^{\rho} \dot{x}^{\nu} - \frac{\dot{e}}{e}\, \dot{x}^{\sigma} = 0, \tag{65}\]

where the Christoffel symbols are still given by the explicit
expressions (53)--(55). The residual reparametrization freedom is
identical, the gauge \(e = \mathrm{const}\) is still available, and in
that gauge (65) collapses to the affine null-geodesic equation

\[\frac{D \dot{x}^{\sigma}}{d\tau} = 0, \qquad g_{\mu\nu}\, \dot{x}^{\mu} \dot{x}^{\nu} = 0. \tag{66}\]

The massless limit is therefore completely regular: the action with
Lagrangian (62) is well-defined, the constraint (64) is enforced
dynamically by the variation of the einbein, and the resulting
trajectories are exactly the null geodesics.

\hypertarget{what-the-null-constraint-does-and-does-not-fix-in-2-1-dimensions}{%
\subsubsection{9.3 What the null constraint does and does not fix in 2 +
1
dimensions}\label{what-the-null-constraint-does-and-does-not-fix-in-2-1-dimensions}}

Here the dimension begins to matter. Write the null constraint (64) in a
locally inertial frame at a point, where
\(g_{\mu\nu} = \eta_{\mu\nu} = \mathrm{diag}(-1, +1, +1)\):

\[- \left( \dot{x}^{0} \right)^{2} + \left( \dot{x}^{1} \right)^{2} + \left( \dot{x}^{2} \right)^{2} = 0. \tag{67}\]

For a future-directed world-line \(\dot{x}^{0} > 0\), and we may
parametrize all solutions by an angle \(\theta \in [0, 2\pi)\):

\[\dot{x}^{1} = \dot{x}^{0} \cos\theta, \qquad \dot{x}^{2} = \dot{x}^{0} \sin\theta. \tag{68}\]

The null constraint is \textbf{one} equation for the \textbf{two}
independent ratios of the three velocity components (the overall scale
is fixed by the choice of parameter, or equivalently by \(e\)). It
therefore leaves a continuous one-parameter family of null directions:
the circle \(S^{1}\) of equation (68). In 1 + 1 dimensions the analogous
computation gives \(\dot{x}^{1} = \pm \dot{x}^{0}\): a discrete set of
two points (formally, the zero-dimensional sphere \(S^{0}\)). In \(n\)
spacetime dimensions the space of null directions is the sphere
\(S^{n-2}\); this counting is developed systematically in Section 10.

The practical consequence for the initial-value problem: in 1 + 1
dimensions the null constraint plus the choice of chirality determine
the velocity up to scale, so the geodesic equation closes on a single
unknown. In 2 + 1 dimensions the direction \(\theta_{0}\) of the initial
null vector is genuine continuous initial data that the equations of
motion cannot supply.

\hypertarget{where-the-dimension-enters-null-directions-mass-shells-and-little-groups}{%
\subsection{10 Where the Dimension Enters: Null Directions, Mass Shells,
and Little
Groups}\label{where-the-dimension-enters-null-directions-mass-shells-and-little-groups}}

This section collects the structural facts that distinguish the 2 +
1-dimensional massless limit from its 1 + 1-dimensional counterpart. The
einbein mechanism is dimension-blind; the kinematics of the constraints
is not. Everything in this section is standard and is cited to the
literature; the derivations we give are elementary.

\hypertarget{the-mass-shell-as-a-homogeneous-space}{%
\subsubsection{10.1 The mass shell as a homogeneous
space}\label{the-mass-shell-as-a-homogeneous-space}}

Work in flat 2 + 1-dimensional Minkowski space with the contravariant
momentum \(p^{\mu} = (p^{0}, p^{1}, p^{2})\) --- the positive-energy
vector on which the Lorentz matrices below act --- and mass shell
(derived in Section 11 below, in the canonical covector variables
\(p_{\mu} = g_{\mu\nu} p^{\nu}\))

\[g_{\mu\nu}\, p^{\mu} p^{\nu} = - (mc_{0})^{2}, \qquad \text{i.e. in an inertial frame} \qquad - (p^{0})^{2} + (p^{1})^{2} + (p^{2})^{2} = - (mc_{0})^{2}. \tag{69}\]

For \(m > 0\) the future sheet of this two-sheeted hyperboloid is the
hyperbolic plane \(H^{2}\): the restriction of the Minkowski quadratic
form to the tangent spaces of the shell is positive-definite, which is
precisely the definition of the hyperboloid model of hyperbolic geometry
(see e.g.~J. G. Ratcliffe, \emph{Foundations of Hyperbolic Manifolds},
Springer, 2nd ed.~2006, Chapter 3). The Lorentz group \(SO(2, 1)\) acts
transitively\footnote{A group \(G\) acts transitively on a set \(X\) if
  any point of \(X\) can be carried to any other by some group element.
  The set \(X\) is then a single orbit, and can be identified with the
  coset space \(G / G_{x}\), where \(G_{x}\) is the stabilizer (little
  group) of any chosen point \(x\).} on this sheet: any future-directed
timelike \(p\) can be boosted and rotated to the rest-frame value
\(p^{(0)} = (mc_{0}, 0, 0)\). The subgroup that leaves \(p^{(0)}\) fixed
consists of the spatial rotations: the \textbf{little group}\footnote{The
  little group (or stabilizer) of a particle's momentum \(p\) is the
  subgroup of the Lorentz group that leaves \(p\) unchanged. Wigner's
  classification (1939) identifies the internal states of an elementary
  particle with the irreducible unitary representations of the little
  group of its momentum.} of a massive particle in 2 + 1 dimensions is

\[G_{\mathrm{massive}} = SO(2). \tag{70}\]

For \(m = 0\) the shell degenerates to the future null cone with the
origin removed. The Lorentz group still acts transitively on it (any
future-null vector can be brought to the standard form
\(k^{(0)} = (\kappa, \kappa, 0)\), \(\kappa > 0\)), but the stabilizer
changes character. The Lorentz transformations that leave \(k^{(0)}\)
fixed form a one-parameter group, generated by the null rotation

\[M = \begin{pmatrix} 0 & 0 & 1 \\ 0 & 0 & 1 \\ 1 & -1 & 0 \end{pmatrix}, \qquad e^{tM} k^{(0)} = k^{(0)} \;\; \forall\, t \in \mathbb{R}, \tag{71}\]

as one checks by direct exponentiation. First,
\(M k^{(0)} = (0,\; 0,\; \kappa - \kappa)^{T} = 0\), so the whole series
collapses to the identity on \(k^{(0)}\). Second, \(M\) is
Lie-algebra-valued for \(SO(2,1)\): with
\(\eta = \mathrm{diag}(-1,+1,+1)\) the condition for a generator is that
\(\eta M\) be antisymmetric, and indeed \(\eta M\) has the nonzero
entries \((\eta M)_{02} = -1 = -(\eta M)_{20}\) and
\((\eta M)_{12} = +1 = -(\eta M)_{21}\). (In the standard basis of
rotation \(J\) in the \(1\)-\(2\) plane and boost \(K_{2}\) along
\(x^{2}\), \(M = K_{2} + J\): the boost and the rotation act oppositely
on \(k^{(0)}\) and cancel.) Third, the action on a general momentum is a
translation-like shear of the transverse coordinate driven by the null
combination: \((e^{tM} p)^{2} = p^{2} + t\,(p^{0} - p^{1})\) ---
exactly, with no higher-order terms, because \(M^{3} = 0\) and
\((M^{2} p)^{2} = 0\) (both verified symbolically in the companion
suite) --- which vanishes exactly on the rays with \(p^{0} = p^{1}\). A
non-compact action rather than a rotation. The stabilizer is therefore
the non-compact one-parameter group

\[G_{\mathrm{massless}} = \mathbb{R} \; (\cong ISO(1)), \tag{72}\]

the degenerate case \(ISO(n-2)\) with \(n = 3\) of the general result
that the little group of a nonzero null vector in \(n\)-dimensional
Minkowski space is the Euclidean group \(ISO(n-2)\) (E. Wigner, ``On
Unitary Representations of the Inhomogeneous Lorentz Group'',
\emph{Annals of Mathematics} \textbf{40} (1939) 149--204, DOI:
10.2307/1968551; a modern textbook account is S. Weinberg, \emph{The
Quantum Theory of Fields}, Vol. 1, Cambridge University Press, 1995,
§2.5).

\hypertarget{spin-in-2-1-dimensions-anyons-for-massive-nothing-for-massless-in-the-conventional-sector}{%
\subsubsection{10.2 Spin in 2 + 1 dimensions: anyons for massive,
nothing for massless in the conventional
sector}\label{spin-in-2-1-dimensions-anyons-for-massive-nothing-for-massless-in-the-conventional-sector}}

The representation theory of the two little groups is completely
different.

\textbf{Massive.} \(SO(2)\) is abelian. Its single-valued irreducible
unitary representations are one-dimensional and labeled by an integer
(\(R(\varphi) = e^{i s \varphi}\) with \(s \in \mathbb{Z}\), since
\(R(\varphi + 2\pi) = R(\varphi)\)); but the object that acts on the
internal state of a relativistic particle is the universal cover of the
Lorentz group, whose rotation subgroup is the universal cover
\(\mathbb{R}\) of \(SO(2)\), and its irreducible unitary representations
are one-dimensional, labeled by any real number \(s\):

\[R(\varphi) \longmapsto e^{i s \varphi}, \qquad \varphi \in \mathbb{R}, \quad s \in \mathbb{R}. \tag{73}\]

In 3 + 1 dimensions single-valuedness for the rotation group (or rather
its double cover) restricts the corresponding label to integers and
half-integers. In 2 + 1 dimensions the fundamental group of the Lorentz
group is \(\pi_{1}(SO(2,1)) = \mathbb{Z}\) --- as opposed to
\(\mathbb{Z}_{2}\) for \(SO(3,1)\) (both statements follow from the
standard deformation retract of \(SO(n,1)\) onto its maximal compact
subgroup \(SO(n)\), with \(\pi_{1}(SO(2)) = \mathbb{Z}\) and
\(\pi_{1}(SO(3)) = \mathbb{Z}_{2}\) tabulated in M. Nakahara,
\emph{Geometry, Topology and Physics}, IOP Publishing, 2nd ed.~2003,
§4.7, Table 4.1) --- so the universal cover imposes no quantization at
all: a massive particle can carry \textbf{anyonic spin}, any real \(s\)
(F. Wilczek, ``Quantum Mechanics of Fractional-Spin Particles'',
\emph{Physical Review Letters} \textbf{49} (1982) 957--959, DOI:
10.1103/PhysRevLett.49.957; relativistic wave equations for anyonic
particles are treated in R. Jackiw and V. P. Nair, ``Relativistic wave
equation for anyons'', \emph{Physical Review D} \textbf{43} (1991)
1933--1942, DOI: 10.1103/PhysRevD.43.1933).

\textbf{Massless.} The little group \(\mathbb{R}\) is non-compact; its
unitary irreducible representations are all one-dimensional
(\(t \mapsto e^{i\rho t}\), \(\rho \in \mathbb{R}\)). A nonzero \(\rho\)
is mathematically allowed: it labels a ``continuous-spin''-type internal
degree of freedom with no localization on any compact direction --- the
2 + 1-dimensional analogue of the continuous-spin representations of the
3 + 1-dimensional Poincaré group --- and explicit 2 + 1-dimensional
field realizations of these nontrivial representations exist (Binegar,
loc. cit. below). The conventional physical choice, going back to
Wigner's 1939 treatment of the 3 + 1-dimensional case, restricts
attention to the \textbf{trivial-translation (finite-spin) sector} in
which the translation-like part of the massless little group acts
trivially on physical states (\(\rho = 0\)); nothing in the
representation theory itself rules the alternative out. In 2 + 1
dimensions, unlike 3 + 1, there is no residual compact rotation factor
left in the little group at all (the group \(SO(1)\) is trivial), so in
this conventional sector a massless particle carries \textbf{no internal
degree of freedom}: its state is unique once its momentum is fixed.

The little-group analysis of relativistic systems in 2 + 1 dimensions is
treated systematically in B. Binegar, ``Relativistic field theories in
three dimensions'', \emph{Journal of Mathematical Physics} \textbf{23}
(1982) 1511--1517, DOI: 10.1063/1.525524.

\hypertarget{the-smooth-classical-limit-and-the-non-smooth-state-counting}{%
\subsubsection{10.3 The smooth classical limit and the non-smooth state
counting}\label{the-smooth-classical-limit-and-the-non-smooth-state-counting}}

We can now state precisely in what sense the massless limit is smooth,
and in what sense it is not --- the question that the einbein
construction was invented to answer.

\begin{itemize}
\tightlist
\item
  \textbf{Classically smooth.} The einbein action (10) depends
  analytically on \(m\) through the single term
  \(- \frac{e}{2}(mc_{0})^{2}\). At \(m = 0\) nothing degenerates: the
  variational principle, the constraint, and the geodesic equation all
  persist, and the massive trajectories approach null geodesics as
  \(m \to 0\) with the initial velocity pushed onto the cone. This is
  the content of Sections 5--9, and it is dimension-independent.
\item
  \textbf{Kinematically continuous, with a change of type.} The space of
  null directions \(S^{1}\) is the boundary of the space of timelike
  directions: the mass shell \(H^{2}\) has the circle as its ideal
  (conformal) boundary at infinity in the hyperboloid model (Ratcliffe,
  loc. cit.). Every null direction is a limit of timelike directions and
  vice versa; no direction is created or destroyed in the limit. What
  changes is the \emph{stabilizer type}: \(SO(2)\) for every \(m > 0\),
  \(\mathbb{R}\) at \(m = 0\) exactly. Group actions are not required to
  vary continuously in their stabilizer type, and here they do not.
\item
  \textbf{Quantum-mechanically non-smooth.} Because the internal states
  are representations of the little group, the state counting jumps at
  \(m = 0\): a massive particle may carry any real anyonic spin \(s\),
  while a massless particle in the conventional trivial-translation
  sector of Section 10.2 carries nothing. This is not a defect of the
  massless theory; it is Wigner kinematics. In 3 + 1 dimensions the same
  phenomenon appears as the jump from \(2s + 1\) massive spin states to
  at most two massless helicity states --- the subject of the companion
  3 + 1 tutorial.
\end{itemize}

\hypertarget{a-remark-on-2-1-dimensional-gravity}{%
\subsubsection{10.4 A remark on 2 + 1-dimensional
gravity}\label{a-remark-on-2-1-dimensional-gravity}}

In three spacetime dimensions the Weyl tensor vanishes identically: in
\(n\) dimensions the Riemann tensor has \(n^{2}(n^{2}-1)/12\)
independent components --- six for \(n = 3\) --- exactly as many as the
Ricci tensor, so the Riemann tensor is algebraically determined by the
Ricci tensor and carries no propagating (traceless) part (see e.g.~S. M.
Carroll, \emph{Spacetime and Geometry}, Chapter 3, for the component
counting). A consequence: with vanishing cosmological constant
(\(\Lambda = 0\)) every vacuum solution of the Einstein equations in 2 +
1 dimensions is locally flat (with \(\Lambda \neq 0\) the vacuum
solutions are instead spacetimes of constant curvature, locally de
Sitter or anti-de Sitter). Independently of \(\Lambda\), pure 2 +
1-dimensional Einstein gravity has no local, propagating graviton
degrees of freedom; in the \(\Lambda = 0\) theory all gravitational
effects are global --- conical defects around point masses and the
scattering of geodesics off them (S. Deser, R. Jackiw, G. 't Hooft,
``Three-dimensional Einstein gravity: dynamics of flat space'',
\emph{Annals of Physics} \textbf{152} (1984) 220--235, DOI:
10.1016/0003-4916(84)90085-X). The einbein formalism of these notes
applies to such backgrounds without modification; the null geodesics of
a conical defect are a clean illustration of Section 9.3's continuous
direction data.

\hypertarget{brief-hamiltonian-perspective}{%
\subsection{11 Brief Hamiltonian
Perspective}\label{brief-hamiltonian-perspective}}

For completeness we record the Hamiltonian formulation that follows from
the einbein Lagrangian. The canonical momenta conjugate to the
coordinates are

\[p_{\lambda} = \frac{\partial L}{\partial \dot{x}^{\lambda}} = \frac{1}{e}\, g_{\lambda\nu}\, \dot{x}^{\nu}. \tag{74}\]

This canonical momentum is a covector. It is related to the
contravariant momentum \(p^{\mu} = \frac{1}{e}\, \dot{x}^{\mu}\) used in
Section 10.1 by index lowering, \(p_{\mu} = g_{\mu\nu}\, p^{\nu}\); with
the signature \((-, +, +)\) this means that in an inertial frame
\(p_{0} = -p^{0}\) and \(p_{i} = p^{i}\), so the component \(p_{0}\) of
the canonical covector is \emph{negative} for a positive-energy
particle. The two index placements give one and the same mass shell,
since \(g^{\mu\nu} p_{\mu} p_{\nu} = g_{\mu\nu} p^{\mu} p^{\nu}\); the
matrices of Section 10.1 act on the contravariant representative
\(p^{\mu}\).

Inverting for the velocities,
\(\dot{x}^{\nu} = e\, g^{\nu\lambda} p_{\lambda}\). The Hamiltonian
(after a Legendre transform) is pure constraint:

\[H = \frac{e}{2} \left( g^{\mu\nu} p_{\mu} p_{\nu} + (mc_{0})^{2} \right). \tag{75}\]

The einbein itself acts as a Lagrange multiplier enforcing the
mass-shell condition

\[g^{\mu\nu} p_{\mu} p_{\nu} + (mc_{0})^{2} = 0. \tag{76}\]

In the massless case this reduces to the light-like condition
\(p^{2} = 0\). In 2 + 1 dimensions the massive shell is the hyperboloid
\(H^{2}\) of Section 10.1 and the massless shell is the punctured null
cone; both are two-dimensional surfaces in the three-dimensional
momentum space, so the constraint surface keeps its dimension across the
limit while changing its geometry (and, with the origin excluded, its
causal character). This constrained Hamiltonian system is the starting
point for the canonical quantization of the relativistic particle; the
little-group representation theory of Section 10 is precisely the
quantum theory of the degrees of freedom that remain after the
constraint (76) is imposed.

\hypertarget{resolution-of-the-equations-of-motion-and-the-resulting-trajectories}{%
\subsection{12 Resolution of the Equations of Motion and the Resulting
Trajectories}\label{resolution-of-the-equations-of-motion-and-the-resulting-trajectories}}

After the residual reparametrization freedom has been used to set
\(e = \mathrm{const}\), the dynamical content of the theory is
completely determined. We summarise the final system of equations and
the nature of its solutions for both the massive and the massless
particle.

\hypertarget{massive-particle-m-0}{%
\subsubsection{\texorpdfstring{12.1 Massive particle
(\(m > 0\))}{12.1 Massive particle (m \textgreater{} 0)}}\label{massive-particle-m-0}}

In the gauge \(e = \mathrm{const}\) the Euler--Lagrange equations reduce
to the pair

\[\frac{D \dot{x}^{\sigma}}{d\tau} = \ddot{x}^{\sigma} + \Gamma^{\sigma}{}_{\rho\nu}\, \dot{x}^{\rho} \dot{x}^{\nu} = 0, \tag{77}\]

\[g_{\mu\nu}\, \dot{x}^{\mu} \dot{x}^{\nu} = - e^{2}(mc_{0})^{2} = \mathrm{const} < 0. \tag{78}\]

Equation (77) is the affine geodesic equation; equation (78) is the
mass-shell constraint.

\textbf{Well-posedness.} Given data \(x^{\mu}(0) = x^{\mu}_{0}\),
\(\dot{x}^{\mu}(0) = u^{\mu}_{0}\) satisfying the mass-shell condition
at \(\tau = 0\), the Picard--Lindelöf theorem guarantees a unique local
solution; the constraint is preserved by (77) because the Levi-Civita
connection is metric-compatible:
\(\frac{d}{d\tau}\left(g_{\mu\nu}\dot{x}^{\mu}\dot{x}^{\nu}\right) = 2\, g_{\mu\nu}\, \dot{x}^{\mu} \frac{D\dot{x}^{\nu}}{d\tau} = 0\).
The solution extends to a maximal interval and is global if the
spacetime is timelike-geodesically complete.

\textbf{Difference from 1 + 1 dimensions.} In 1 + 1 dimensions the
constraint (78) restricts the velocity to a one-parameter family ---
explicitly
\(\dot{x}^{1} = \pm \sqrt{(\dot{x}^{0})^{2} - e^{2}(mc_{0})^{2}}\),
equivalently a single continuous rapidity --- so \textbf{one} continuous
shell parameter remains. In 2 + 1 dimensions the shell is
two-dimensional: the constraint fixes only the overall scale of the
tangent vector; the \textbf{two} independent ratios of the three
velocity components --- the direction angle \(\theta_{0}\) in the
\((\dot{x}^{1}, \dot{x}^{2})\) plane and the speed ratio --- are free
initial data, and (77) remains a genuine system of coupled second-order
equations that the constraint alone cannot reduce to a single
first-order ODE. If the metric admits a Killing vector field\footnote{A
  Killing vector field is a vector field \(\xi^{\mu}\) along whose flow
  the metric is invariant, equivalently
  \(\nabla_{(\mu} \xi_{\nu)} = 0\). The inner product of a Killing
  vector with the tangent vector of a geodesic is constant along the
  geodesic; see R. M. Wald, \emph{General Relativity} (University of
  Chicago Press, 1984), Appendix C, or Carroll, \emph{Spacetime and
  Geometry}, Appendix B.} \(\xi^{\mu}\), the quantity
\(E = g_{\mu\nu} \xi^{\mu} \dot{x}^{\nu}\) is conserved and supplies one
further independent relation; each additional independent commuting
Killing vector supplies another first integral of the same form.

\textbf{Flat-space illustration.} When \(g_{\mu\nu} = \eta_{\mu\nu}\)
the Christoffel symbols vanish and the equations become
\(\ddot{x}^{\sigma} = 0\) with
\(\eta_{\mu\nu} \dot{x}^{\mu} \dot{x}^{\nu} = -e^{2}(mc_{0})^{2}\). The
general solution is the straight line

\[x^{\sigma}(\tau) = x^{\sigma}_{0} + u^{\sigma}_{0}\, \tau, \qquad \eta_{\mu\nu} u^{\mu}_{0} u^{\nu}_{0} = - e^{2} (mc_{0})^{2}, \tag{79}\]

with the free data \((x^{\mu}_{0}, u^{\mu}_{0})\) including the
direction angle \(\theta_{0}\) of Section 9.3. This is the exponential
map of a vector space,
\(x(\tau) = \exp_{x_{0}}(\tau u_{0}) = x_{0} + \tau u_{0}\).

\hypertarget{massless-particle-m-0}{%
\subsubsection{\texorpdfstring{12.2 Massless particle
(\(m = 0\))}{12.2 Massless particle (m = 0)}}\label{massless-particle-m-0}}

In the same gauge the equations reduce to

\[\frac{D \dot{x}^{\sigma}}{d\tau} = 0, \qquad g_{\mu\nu}\, \dot{x}^{\mu} \dot{x}^{\nu} = 0. \tag{80}\]

Any solution is a null geodesic of the background metric, affinely
parametrized. The initial-value problem is well-posed given data
\((x^{\mu}_{0}, k^{\mu}_{0})\) with \(k_{0}\) future-null; the direction
of \(k_{0}\) --- a point of the circle \(S^{1}\) of null directions ---
is part of the data, not fixed by any constraint. In Minkowski space the
solutions are the straight null lines

\[x^{\sigma}(\tau) = x^{\sigma}_{0} + k^{\sigma}_{0}\, \tau, \qquad \eta_{\mu\nu} k^{\mu}_{0} k^{\nu}_{0} = 0, \qquad k^{0}_{0} > 0. \tag{81}\]

A useful conformal fact, true in every dimension but first visible as a
nontrivial statement here: the \emph{unparametrized} null geodesics of a
metric depend only on its conformal class. Under a Weyl rescaling
\(g_{\mu\nu} \to \Omega^{2}(x)\, g_{\mu\nu}\) the null cone is unchanged
pointwise, and the null geodesic equation is preserved up to a change of
parameter (R. M. Wald, \emph{General Relativity}, Appendix D). The
einbein makes this transparent at the level of the action: under the
combined transformation

\[g_{\mu\nu} \longmapsto \Omega^{2}(x)\, g_{\mu\nu}, \qquad e \longmapsto \Omega^{2}(x(\tau))\, e(\tau) \tag{82}\]

--- the second being a Weyl rescaling of the world-line metric (9) by
the pullback of the same factor --- the massless Lagrangian (62) is
invariant pointwise. For \(m > 0\) no such invariance exists: the mass
term \(-\frac{e}{2}(mc_{0})^{2}\) picks up a factor \(\Omega^{2}\). This
is the action-level shadow of the fact that timelike geodesics are not
conformally invariant while unparametrized null geodesics are.

\hypertarget{summary-of-the-trajectories}{%
\subsubsection{12.3 Summary of the
trajectories}\label{summary-of-the-trajectories}}

\begin{longtable}[]{@{}
  >{\raggedright\arraybackslash}p{(\columnwidth - 6\tabcolsep) * \real{0.2500}}
  >{\raggedright\arraybackslash}p{(\columnwidth - 6\tabcolsep) * \real{0.2500}}
  >{\raggedright\arraybackslash}p{(\columnwidth - 6\tabcolsep) * \real{0.2500}}
  >{\raggedright\arraybackslash}p{(\columnwidth - 6\tabcolsep) * \real{0.2500}}@{}}
\toprule\noalign{}
\begin{minipage}[b]{\linewidth}\raggedright
Case
\end{minipage} & \begin{minipage}[b]{\linewidth}\raggedright
Final equations
\end{minipage} & \begin{minipage}[b]{\linewidth}\raggedright
Geometric character
\end{minipage} & \begin{minipage}[b]{\linewidth}\raggedright
Free direction data
\end{minipage} \\
\midrule\noalign{}
\endhead
\bottomrule\noalign{}
\endlastfoot
Massive (\(m > 0\)) & affine geodesic eq. + mass shell & future-directed
timelike geodesic & angle \(\theta_{0} \in S^{1}\) \\
Massless (\(m = 0\)) & affine geodesic eq. + null condition & null
geodesic & angle \(\theta_{0} \in S^{1}\) \\
\end{longtable}

In both cases the einbein formulation has achieved its purpose: the
equations of motion are well-defined, the massless limit is continuous,
and the solutions are precisely the geodesics of the appropriate causal
type. What 2 + 1 dimensions add to the 1 + 1 story is the continuous
circle of directions and the little-group analysis of Section 10.

\hypertarget{summary}{%
\subsection{13 Summary}\label{summary}}

We have shown, by completely explicit algebraic manipulations, that the
einbein reformulation of the relativistic point particle extends from 1
+ 1 to 2 + 1 dimensions without any modification of its variational
core: the auxiliary-field equation (25), the recovery of the square-root
action (Section 7), the geodesic equation with its non-affine term (56),
and the gauge \(e = \mathrm{const}\) that removes it are all
dimension-independent. The new content of the 2 + 1-dimensional theory
is kinematic and group-theoretic:

\begin{itemize}
\tightlist
\item
  the null directions form a circle \(S^{1}\) (Section 9.3), so the
  direction of a massless particle is continuous initial data rather
  than a discrete chirality;
\item
  the massive mass shell is the hyperbolic plane \(H^{2}\) with little
  group \(SO(2)\), whose representations allow any real spin --- anyons
  (Section 10.2);
\item
  the massless little group is the non-compact group \(\mathbb{R}\),
  whose nontrivial unitary representations are mathematically allowed;
  in the conventional trivial-translation sector, physical massless
  particles carry no spin degree of freedom (Section 10.2);
\item
  the massless limit is smooth at the level of the action and the
  geodesics, while the internal-state counting jumps --- a
  Wigner-kinematic fact, not a dynamical pathology (Section 10.3);
\item
  the massless einbein action is Weyl-invariant, exposing the conformal
  nature of unparametrized null geodesics (Section 12.2).
\end{itemize}

The 3 + 1-dimensional case --- in which the little groups become
\(SO(3)\) and \(ISO(2)\), the null directions form the celestial sphere
\(S^{2}\), and the helicity structure of massless particles appears ---
is treated in the companion tutorial.

\hypertarget{references}{%
\subsection{References}\label{references}}

{[}0{]} T. Fulceri, ``Transforming the Canonical Lagrangian of a
Relativistic Particle for a Smooth Massless Limit: The Einbein
Formulation in 1 + 1 Dimensions'', Zenodo technical tutorial notes, 11
August 2026. DOI: 10.5281/zenodo.21879560.

{[}1{]} J. Polchinski, \emph{String Theory, Volume 1: An Introduction to
the Bosonic String}, Cambridge University Press, Cambridge, 1998.
(Chapter 1: the einbein formulation as preparation for the Polyakov
action.) DOI: 10.1017/CBO9780511816079.

{[}2{]} B. Zwiebach, \emph{A First Course in String Theory}, Cambridge
University Press, Cambridge, 2nd ed.~2009. (Chapter 5: relativistic
point-particle action and reparametrization invariance.) DOI:
10.1017/CBO9780511841682.

{[}3{]} E. Wigner, ``On Unitary Representations of the Inhomogeneous
Lorentz Group'', \emph{Annals of Mathematics} \textbf{40} (1939)
149--204. DOI: 10.2307/1968551.

{[}4{]} S. Weinberg, \emph{The Quantum Theory of Fields, Volume 1:
Foundations}, Cambridge University Press, 1995. (Chapter 2: little
groups; §2.5: massless particles.) DOI: 10.1017/CBO9781139644167.

{[}5{]} F. Wilczek, ``Quantum Mechanics of Fractional-Spin Particles'',
\emph{Physical Review Letters} \textbf{49} (1982) 957--959. DOI:
10.1103/PhysRevLett.49.957.

{[}6{]} R. Jackiw and V. P. Nair, ``Relativistic wave equation for
anyons'', \emph{Physical Review D} \textbf{43} (1991) 1933--1942. DOI:
10.1103/PhysRevD.43.1933.

{[}7{]} B. Binegar, ``Relativistic field theories in three dimensions'',
\emph{Journal of Mathematical Physics} \textbf{23} (1982) 1511--1517.
DOI: 10.1063/1.525524.

{[}8{]} S. Deser, R. Jackiw, G. 't Hooft, ``Three-dimensional Einstein
gravity: dynamics of flat space'', \emph{Annals of Physics} \textbf{152}
(1984) 220--235. DOI: 10.1016/0003-4916(84)90085-X.

{[}9{]} R. M. Wald, \emph{General Relativity}, University of Chicago
Press, 1984. (Chapter 3: curvature; Appendix C: Killing vectors;
Appendix D: conformal transformations and null geodesics.) ISBN:
9780226870335.

{[}10{]} S. M. Carroll, \emph{Spacetime and Geometry: An Introduction to
General Relativity}, Cambridge University Press, 2019. (Chapter 3:
geodesics, Christoffel symbols, curvature component counting.) ISBN:
9781108488396.

{[}11{]} J. G. Ratcliffe, \emph{Foundations of Hyperbolic Manifolds},
Springer, 2nd ed.~2006. (Chapter 3: the hyperboloid model and its ideal
boundary.) DOI: 10.1007/978-0-387-47322-2.

{[}12{]} F. Bastianelli, \emph{Relativistic Quantum Mechanics and Path
Integrals}, lecture notes, University of Bologna, 2023--2024.
https://www-th.bo.infn.it/people/bastianelli/2-Mechanics-RQMPI-23-24.pdf
(accessed August 2026). (Explicit treatment of the einbein action.)

{[}13{]} M. Nakahara, \emph{Geometry, Topology and Physics}, IOP
Publishing, 2nd ed.~2003. (§4.7, Table 4.1: fundamental groups of the
classical groups.) ISBN: 9780750306065.

\emph{End of tutorial notes.}

\end{document}
